\chapter{Тренажеры автотранспорта и спецтехники с имитацией движения. Необходимость создания систем обратной связи} \label{chapt1}
Современные тренажеры представляют собой программно-аппаратные комплексы (ПАК), используемые для обучения персонала управлению сложными техническими системами и объектами. Использование тренажеров эффективно в тех случаях, когда ошибка обучаемого может повлечь серьезные повреждения техники, а также для экономии временных и энергоресурсов, связанных с подготовкой и проведением обучения на реальной технике. 
Современные программные и аппаратные средства тренажеров также позволяют осуществлять имитацию различных нештатных ситуаций (отказ оборудования, неблагоприятные погодные условия), что на реальной технике в принципе нереализуемо.

В тренажерах автотранспорта и спецтехники (АТ и СТ), как в одном из наиболее востребованных в наши дни направлений разработки тренажеров, широко применяются средства тактильной обратной связи (ТОС) с обучаемым. ТОС включает как простую ОС, реализуемую с помощью вибромоторов (т.н. Force Feedback), так и полномасштабную имитацию акселерационных эффектов движения с помощью систем подвижности (СП). 

Тренажеры АТ и СТ с точки зрения возможностей по имитации движений подразделяются на два класса - статические и динамические. Динамические тренажеры отличаются наличием системы подвижности, имеющей от одной до шести степеней свободы. В некоторых специфических конфигурациях тренажеров могут применяться СП с большим числом степеней свободы, часть из которых при этом являются независимыми (например, фирма MOOG предлагает вариант компоновки СП MB-E-6DOF собственного производства на подвижном рельсовом основании для использования в автотренажерах).

Независимо от варианта реализации, динамические тренажеры выгодно отличаются:
\begin{itemize}
\item{более реалистичными ощущениями при управлении техникой};
\item{большей методической ценностью засчет формирования необходимых моторных навыков у обучаемого};
\item{могут быть использованы для повышения квалификации операторов и переподготовки специалистов}.
\end{itemize}

К недостаткам динамических тренажеров относятся:
\begin{itemize}
\item{высокая стоимость};
\item{повышенные габариты, необходимость специального помещения для размещения тренажера};
\item{повышенная потребляемая мощность};
\item{сложность эксплуатации, необходимость обслуживания и планового ремонта};
\item{сложность транспортировки и инсталляции}.
\end{itemize}

Несмотря на вышеперечисленные недостатки, тренажеры с системами подвижности пользуются большим спросом на рынке тренажерных систем, вследствие чего ведущие производители тренажеров АТ и СТ стремятся оптимизировать конструкцию и энергопотребление СП, уменьшить эксплуатационные расходы, упростить инсталляционные и сервисные процедуры.

В настоящее время существующие на рынке конфигурации СП можно классифицировать в соответствии со следующей схемой:
\begin{figure}[!h]
	\center
	\includegraphics[width=1.0\textwidth]{MPclassification.png}
	\caption{Классификация основных типов СП}
	\label{img:MPclassification}
\end{figure}

Рассмотрим подробнее типы механизмов и приводов СП.

\section{Разновидности механизмов СП, применяемых в тренажерах АТ и СТ} \label{sect1_1}
Одной из основных характеристик динамической платформы является количество степеней свободы, которое для механизмов параллельной структуры, определяется количеством приводных звеньев.

Динамическая платформа, изображенная на рисунке \ref{img:3dofKT} обеспечивает 3 степени свободы (крен, тангаж, подъем) и позволяет отработать большую часть необходимых эффектов для целей обучения операторов техники. Приводы СП через кривошипно-шатунные механизмы осуществляют поступательное движение трех точек подвижной платформы, при этом геометрический центр конструкции соединен с основанием подпружиненной опорой, разгружающей двигатели. Недостатком данной схемы является наличие паразитных колебаний, которые в определенных условиях могут возникать при отработке интенсивных вертикальных движений и тангажировании.

\begin{figure}[ht]
    \center{
        \hfill
        \subcaptionbox[List-of-Figures entry]{\label{img:3dofKT}} {\includegraphics[width=0.45\linewidth]{3dofKT}}%
        \hfill       
        \subcaptionbox{\label{img:2dofKT}} {\includegraphics[width=0.45\linewidth]{2dofKT}}
        \hfill
    }

    Система подвижности с двумя (а) и тремя (б) степенями свободы производства <<Кронштадт Технологии>>.
    \caption{СП 2DOF и 3DOF } % Этот текст попадает в названия рисунков в списке рисунков
    \label{img:3dof2dofKT}
\end{figure}

Упрощенным вариантом первой схемы является вариант, изображенный на рис \ref{img:2dofKT}, обеспечивающий 2 степени свободы (крен, тангаж). Такая конструкция является более жесткой за счет закрепления верхнего основания на опоре с карданным шарниром. Использование двух двигателей вместо трех делает СП более экономичной. Однако в такой схеме горизонтальное положение подвижной платформы является неустойчивым, т.к. при отсутствии постоянного удерживающего момента в момент постановки на тормоза кривошипы перемещаются в нижнюю мертвую точку под действием силы тяжести. Данная особенность затрудняет безударную постановку СП на тормоза.

Дальнейшим упрощением двухстепенной СП является динамическая платформа с одним приводом вращательного типа, изображенная на рисунке \ref{img:1dofEOV}. Данный тип СП предназначен для использования в компактных тренажерах спецтехники и позволяет имитировать колебания по единственной оси -- тангаж вперед/назад.

\begin{figure}
\center
\includegraphics[width=\linewidth]{excavator1DOF}
\caption{Одностепенная СП тренажера одноковшового экскаватора}
\label{img:1dofEOV}
\end{figure}

\begin{figure}[ht]
    \center{
        \hfill
        \subcaptionbox[List-of-Figures entry]{\label{img:4dofKT}} {\includegraphics[width=0.45\linewidth]{4dofKT}}%
        \hfill       
        \subcaptionbox{\label{img:4dof_murom}} {\includegraphics[width=0.45\linewidth]{4dof_murom}}
        \hfill
    }

    Система подвижности 3DOF с дополнительным независимым приводом
    \caption{СП 4DOF} % Этот текст попадает в названия рисунков в списке рисунков
    \label{img:4dofplatforms}
\end{figure}

Для обеспечения дополнительной степени свободы может применяться вспомогательный привод, как это реализовано в системах подвижности производства <<Кронштадт технологии>> (рисунок \ref{img:4dofKT}) и Муромского радиозавода (рисунок \ref{img:4dof_murom}).
Наличие дополнительного вращающегося тела на подвижной платформе влечет за собой сложности в технической реализации конструкции, увеличение габаритов и потребляемой мощности. Однако, данный тип механизма СП нашел применение в экипажных тренажерах автотехники. 

\begin{figure}[ht]
    \center{
        \hfill
        \subcaptionbox[List-of-Figures entry]{\label{img:6dofKT}} {\includegraphics[width=0.45\linewidth]{6dofKT}}%
        \hfill       
        \subcaptionbox{\label{img:6dofMoog}} {\includegraphics[width=0.45\linewidth]{Moog.png}}
        \hfill
    }

    Система подвижности 6DOF с пневмо (а) и электроприводом (б)
    \caption{СП 6DOF} % Этот текст попадает в названия рисунков в списке рисунков
    \label{img:6dofplatforms}
\end{figure}

В целях увеличения спектра имитируемых эффектов в тренажерах также применяются кинематические схемы с шестью степенями свободы (рисунок \ref{img:6dofplatforms}), реализованные на базе линейных или вращательных приводов. Данный тип платформ позволяет имитировать такие специфические движения, как вращение закрепленного на подвижной платформе объекта относительно вертикальной плоскости и его поступательные смещения, что является критичным  при моделировании поворотов и заносов техники.

\section{Разновидности приводных систем СП} \label{sect1_2}
Системы подвижности подразделяются по типу привода на гидравлические, электрические, электрогидравлические и электропневматические.

Гидравлические системы подвижности наиболее широко применяются при производстве авиасимуляторов в связи с их очевидными преимуществами:
\begin{itemize}
\item высокими развиваемыми усилиями, что позволяет достигать больших ускорений и перемещать объекты значительной массы (развиваемые силы достигают десятков тс);
\item высокой плавностью хода;
\item низким износом механической части привода засчет смазывания трущихся частей рабочей жидкостью.  
\end{itemize} 
К недостаткам гидравлических СП относят:
\begin{itemize}
\item высокую стоимость эксплуатации;
\item высокую потребляемая мощность;
\item необходимость выделенного помещения для гидронасосной станции, применения дополнительных мер шумо-, виброизоляции.
\end{itemize}
Гидравлические системы в целом малопрактичны и требуют регулярного обслуживания -- устранения утечек рабочей жидкости, герметизации соединений трубопроводов и т.д. 

Электромеханические СП, напротив:
\begin{itemize}
\item обладают низким энергопотреблением;
\item практичны -- обслуживание включает лишь периодическое смазывание механических частей;
\item компактны. Как правило, шкаф управления может быть размещен в непосредственной близости от платформы или смонтирован на нижнем основании СП.
\end{itemize}

СП с электроприводом получили наибольшее распространение в компактных тренажерах АТ и СТ, развлекательных аттракционах.

Несмотря на то, что электрические СП соответствуют требованиям международных стандартов в области разработки симуляторов, по показателям эффективности имитации движения они ощутимо уступают гидравлическим СП. Приводы электромеханической СП зачастую неспособны развивать достаточную мгновенную мощность, в то время как гидропривод обеспечивает лучшие характеристики при аналогичных габаритах и способен накапливать энергию за счет гидроаккумулятора. Гидравлические СП признаются экспертами эталоном качества движения. 

В данном типе СП были применены решения, позволяющие уменьшить энергопотребление и сделать систему подвижности более компактной засчет размещения индивидуального гидронасоса в непосредственной близости от каждого из гидроцилиндров платформы в отличие от классической гидравлической СП, в которой организована централизованная подача жидкости от одной насосной станции. Каждый из насосов электрогидравличекой СП имеет привод от электромотора, также смонтированного вблизи платформы. Такая компоновка гидравлики позволяет в том числе уменьшить потери давления за счет сокращения длины трубопровода от насосов к гидроцилиндрам. Таким образом, L-3 Communications удалось компенсировать часть недостатков классической гидравлической СП, указанных выше. Плавность хода и качество гидбридной СП оказалось близкими к гидравлической. 

Пневматические системы подвижности до недавнего времени практически не использовались в связи с общеизвестными недостатками, такими как низкий КПД, сложность позиционного управления и высокая стоимость сжатого воздуха. Однако в 2011 
году компания <<Логос>> использовала в своем тренажере вождения <<КАМАЗ>> пневматическую шестистепенную систему подвижности на базе оборудования Festo. Применение пневмопривода позволило добиться более качественной имитации поведения подвески транспортного средства за счет упругих свойств сжатого воздуха. 
 
На рынке также представлена электропневматическая система подвижности фирмы Moog, в которой реализован принцип пневматической разгрузки червячной пары.

\section{Анализ характеристик АТ и СТ с целью формирования требований к приводам СП} \label{sect1_3}

При формировании закона движения платформы учитываются реальные параметры движения техники -- поступательные и угловые перемещения, скорости и ускорения по осям $XYZ$. В связи с ограничениями, накладываемыми СП, производится имитация кратковременной перегрузки с последующим плавным возвратом в исходное положение.

\begin{figure}
\center
\includegraphics[width=0.75\linewidth]{washoutFilterAlgorithm}
\caption{К описанию washout-фильтра}
\label{img:WashoutFilterAlgo}
\end{figure}

Данный принцип проиллюстрирован на рисунке \ref{img:WashoutFilterAlgo}. На участке 1 производится разгон (толчок), близкий к реальному ускорению ТС. На участке 2 показано расхождение движения реальной техники и подвижности тренажера. Подвижность плавно тормозит и перемещается в обратном направлении таким образом, чтобы ускорение в фазе 2-3 было ниже порога восприятия оператора. К моменту времени 3 СП занимает исходную позицию и готова к воспроизведению следующего акселерационного эффекта. <<Триггером>> к последующему началу движения СП может быть какое-либо событие, регистрируемое математической моделью ТС -- взаимодействие с виртуальными объектами, изменение сигналов от  органов управления и т.д.

На рисунке \ref{img:motionCueingAlgorithm} проиллюстрирован классический алгоритм формирования движения СП. Входными сигналами являются поступательные силы, действующие на ТС ($f_{AA}$) и угловые скорости ТС ($\omega_{AA}$), рассчитываемые на стороне математической модели тренажера. После масштабирования (коэффициенты $K_f$ и $K_{\omega}$), результирующие поступательные ускорения раскладываются на высокочастотную и низкочастотную компоненты с помощью фильтров верхних и нижних частот (на схеме -- фильтры ФВЧ и ФНЧ соответственно). Компонент $a_2$, содержащий высокие гармоники поступательных ускорений, проходит через двойной интегратор, с выхода которого снимается поступательное перемещение динамической платформы по осям $X,~Y,~Z$.

Низкочастотный компонент $f_{L}$, используется для расчета углов ориентации платформы $\beta_{L}$ путем проецирования вектора кажущегося ускорения на координатные оси. На выход канала поступает вектор $\beta_{SL}$, ограниченный по скорости нарастания (<<RL>> -- Rate limiter на схеме).

В канале угловых скоростей  производится выборка высших гармоник (ФВЧ) с последующим интегрированием и получением высокочастотной компоненты углов ориентации $\beta_{SH}$, которая суммируется с $\beta_{LH}$. Результат суммирования используется для задания углового положения динамической платформы.
%При разработке математического закона движения платформы тренажера моделируемое на тренажере уравнение движения ВС с помощью методов гармонического анализа раскладывают в ряд гармонических колебаний — гармоник. Первые гармоники — гармоники самой малой частоты вносят наибольший вклад в перемещение самолета. При этом человек именно к этим длиннопериодическим колебаниям наименее чувствителен. Так, если медленно увеличивать перегрузку до небольших значений, то человек в положении сидя может даже не почувствовать её изменения. Высшие гармоники с ростом частоты вносят все меньший вклад в движение и они все более чувствительны для человека. Поэтому низшие гармоники подавляют, используя фильтр верхних частот.

%Динамика движения платформы тренажера продемонстрирована на графике. На графике видно, что система подвижности имитирует перегрузку на небольшом участке времени (меньше секунды), на котором ускорение разгона платформы достигает ускорения моделируемого ВС. Далее в связи с ограниченностью рабочего хода платформы, она тормозится и возвращается в нейтральное положение. При этом, торможение и возврат платформы осуществляется с ускорением ниже порога восприятия человека.

\begin{figure}
	\center
	\includegraphics[width=\linewidth]{motionCueingAlgorithm}
	\caption{Классическая схема формирования движения СП}
    \label{img:motionCueingAlgorithm}
\end{figure}

%Помимо кратковременной имитации перегрузки существует также возможность и долговременной имитации перегрузки. Наиболее простым и широко используемым способом имитации длительной перегрузки является использование горизонтальной составляющей силы тяжести для имитации продольной и боковой перегрузки путём соответствующего наклона платформы. Для того, чтобы добиться этого эффекта, при формирования закона движения платформы уравнение движения ВС пропускают через фильтр нижних частот, который подавляет высшие гармоники.

Для анализа спектра имитируемых эффектов и последующего формирования тестовых траекторий движения СП выделены два наиболее характерных для своего класса тренажера – тренажер оператора экскаватора карьерного гусеничного ЭКГ-18Р и тренажер водителя многоосного автомобиля БТР-82А .
В таблице \ref{tbl:OperatorEffects} приведен перечень основных эффектов, необходимых для воспроизведения на месте оператора в тренажерах ЭКГ-18Р и БТР-82А.

\begin{table} [htbp]%
    \centering
	\caption{Эффекты на месте оператора АТ и СТ}%
	\label{tbl:OperatorEffects}
    \renewcommand{\arraystretch}{1.5}
    \begin{SingleSpace}
	\begin{tabulary}{\textwidth}{@{}>{\zz}L >{\zz}L >{\zz}L >{\zz}L >{\zz}L@{}}
        \toprule     %%% верхняя линейка
    	п/п	& Эффект & Момент воспроизведения (ЭКГ-18Р)& Момент воспроизведения (БТР-82А)\\
        \midrule
    	1 & Равномерные вибрации & при ходе экскаватора & при работе двигателя на холостом ходу \\
    	2 & Вибрации по высоте различной частоты и амплитуды & при копании, добыче породы & при движении по неровной поверхности (различные типы грунта)	\\
    	3 & Резкий толчок & при ударе о грунт или другие предметы окружения & при столкновениях, при наезде на препятствия\\
    	4 & Центростремительное ускорение & при повороте кузова экскаватора & при повороте техники, заносах\\
    	5 & Изменение положения техники в пространстве & при движении по забою (незначительное) & при движении по местности (неровности рельефа) \\
        \bottomrule
	\end{tabulary}%
   	\end{SingleSpace}
\end{table}

На рисунке \ref{img:effectsTypical} показаны перемещения и скорости, рассчитываемые в математической модели БТР-82А 
(имитация проезда канавы - таблица \ref{tbl:OperatorEffects}, п.5) и модели ЭКГ-18Р (воспроизведение удара о грунт - таблица \ref{tbl:OperatorEffects}, п.3).

\begin{figure}[!htbp]
    \center{
        \subcaptionbox[List-of-Figures entry]{\label{img:effectsTypicalAuto}} {\includegraphics[width=0.9\linewidth]{effects_typical_auto}}%
        \hfill       
        \subcaptionbox{\label{img:effectsTypicalSpec}} {\includegraphics[width=0.9\linewidth]{effects_typical_spec}}
    }

    Траектории движения для тренажеров авто- (а) и спецтехники (б). 1 -- заданный крен, 2 -- заданный тангаж платформы
    \caption{Типовые имитационные эффекты тренажеров АТ и СТ} % Этот текст попадает в названия рисунков в списке рисунков
    \label{img:effectsTypical}
\end{figure}

В зависимости от требований к тренажеру и выбранной конфигурации динамической платформы, спектр имитируемых эффектов может быть весьма ограничен в связи с техническими ограничениями. Например, достаточно редко воспроизводится реальная виброактивность двигателя ТС, т.к. данный тип обратной связи требует установки одного или нескольких вибромоторов, что сказывается на стоимости изделия и энергопотреблении. Зачастую в практических применениях разработчики ограничиваются лишь имитацией низкочастотных движений, т.к. электроприводы платформы не позволяют создать достаточное ускорение для полноценной отработки удара о препятствие и подобных эффектов.

\section{Постановка задачи} \label{sect1_4}

Среди рассмотренных вариантов реализации СП для тренажеров АТ и СТ, стоит выделить трехстепенные электромеханические платформы с кривошипно-шатунным приводом звеньев, которые в силу особенностей выбранного привода и кинематической схемы обладают оптимальными характеристиками с точки зрения энергопотребления, затрат на эксплуатацию и обслуживание.  К основным недостаткам данного типа СП относят сравнительно низкое качество имитации движений, требующих больших ускорений – реакция с запозданием при проезде препятствий, неправдоподобные ощущения при ударах модели техники об объекты на 3D-сцене, пр. Данная особенность как правило связана с невысоким значением перегрузочной способности по вращающему моменту в используемых асинхронных двигателях. Помимо этого, при определенной конфигурации кабины тренажера (высоко расположенный центр тяжести, большая масса кабины), мотор-редукторы СП могут длительное время находиться в пусковом режиме работы, что ведет к повышенному потреблению энергии и нестабильности системы. В связи с этим, повышенное внимание при разработке алгоритмов управления СП уделяется точному расчету рабочей зоны как по кинематическим, там и динамическим ограничениям платформы.

\begin{figure}
	\center
	\caption{Обобщенная схема УСП ТК АТ и СТ}
	\label{img:CommonUniversalPlatformControl}
	\includegraphics[width=\linewidth]{control_system_common}
\end{figure}

На рисунке \ref{img:CommonUniversalPlatformControl} представлена типовая структурная схема унифицированной  системы подвижности (УСП) тренажера АТ и СТ. Схема включает:
\begin{itemize}
\item Динамическую платформу подвижности (ДПП);
\item Клиент СП -- часть программного комплекса, реализующая алгоритм, описанным в пункте \ref{sect1_3};
\item ЛСУ ИМ -- локальные системы управления исполнительными механизмами. Реализуются в виде отдельных САУ  приводов платформы; 
\item КД -- контроллер движения. Реализует обработку команд от клиента УСП, ограничение задаваемых перемещений, скоростей и ускорения платформы и совместное управление ЛСУ ИМ.
\end{itemize}

Характеристики качества имитации акселерационных эффектов всецело зависят от алгоритмов управления, реализуемых контроллером движения платформы.
Рассмотренная в \cite{VarkovPhD,vakbib1,vakbib2} классическая схема контроллера движения с независимыми каналами опирается лишь на уравнения обратной или прямой задач кинематики, вследствие чего ее применение влечет вышеперечисленные недостатки. Для их устранения, требуется применение подходов, основанных на более детализированной модели объекта, учитывающей массогабаритные и инерционные характеристики динамической платформы с закрепленным на ней объектом. Подобные исследования проводились в большом количестве работ \cite{Merlet2006, Meng2010, Lukanin2005, Lee2003, He2007}. Ряд авторов используют подход, основанный на расчете движущих моментов и введении в регулятор ошибки по движущим моментам. Часть работ посвящены нейросетевым и нечетким регуляторам \cite{Wang,Singh2014,Salisbury2016}. Данная работа посвящена исследованию одного из вариантов усовершенствованной системы управления, описанном в \cite{Zenkevich2004}, который заключается в коррекции траектории движения в зависимости от ограничений, накладываемых динамикой механической части СП. Используемый метод основывается на нахождении точного решения обратной задачи динамики (ОЗД).  

К системе управления при этом предъявляются следующие требования:
\begin{itemize}
	\item минимальное запаздывание (расхождение между анимацией движения техники и движением платформы должно составлять не более 100-200 мс); 
	\item отсутствие перерегулирования;
	\item устойчивость при отработке задающих воздействий в диапазоне частот до 1 Гц.
\end{itemize}

Общей целью работы, как отмечалось ранее, является разработка унифицированной системы управления (УСП), позволяющей абстрагировать от клиента системы подвижности все детали процесса управления платформой и предоставить единый интерфейс для создания самодостаточного ПАК управления платформой.